\chapter{Cross-Asset Lead--Lag}\label{ch:s1:cross-asset-lead-lag}

The same piece of news reaches different prices at different speeds. Hong, Torous and Valkanov found that the returns of industries such as retail,
metal and petroleum forecast the whole US stock market by up to two months; Hou found that small firms follow the big firms of their own industry, most
slowly after bad news; at the other end of the clock, futures and the stocks or funds that track them are linked by leads of seconds. On this chapter's
synthetic markets a lead of five days from commodities to three industries is found in two of its three links by a careful scan and traded at a Sharpe
ratio of 1.73 after costs out of sample; a lead of two seconds between two assets is found in every session and pays a trader who is at most one second
late. The build is \texttt{firm.xasset}.

\section{Intermarket signals at daily horizons}

\begin{definition}[Intermarket signal, slow diffusion]\label{def:s1:cross-asset-lead-lag:intermarket}
An \emph{intermarket signal}\index{intermarket signal} predicts the returns of one market from the past returns of another that reflects the same
fundamentals sooner, such as a commodity for its producers or a bond market for rate-sensitive stocks. \emph{Slow diffusion}\index{slow diffusion} is the
mechanism behind it: information reaches the prices of different assets at different speeds because investors specialise, attend to few markets, or
face costs of trading, so prices that should move together move in sequence.
\end{definition}

A lead--lag relationship (Book~7, chapter~10) is found by correlating one series with the other's past. At daily horizons the candidates are many and
the true links few. The chapter plants them: \texttt{firm.synthmkt}'s ten industry factor returns follow \texttt{firm.synthfut}'s ten commodities
through three links, each industry receiving 0.08 times the average of one commodity's returns over the previous five days
(\cref{lst:s1:cross-asset-lead-lag:scan}). The effect is small: with commodity volatility of about 25\% and industry volatility of 8\%, each day's lag
carries a correlation of about 0.05, a t-statistic below 2 in five years.

Two scans run on the first five years. The lag-by-lag scan correlates every industry with every commodity at lags of 1 to 10 days, 1\,000 tests, with a
Bonferroni threshold of 4.06 for a 5\% family-wise error. It finds one link, the first planted one, at the right lag of five days ($t = 4.7$); the
other two planted links peak at 2.7 and 3.5, below the threshold (\cref{fig:s1:cross-asset-lead-lag:scan}). The window scan tests one statistic per
pair, the correlation of each industry with the commodity's past five-day return, 100 tests with a threshold of 3.48. It finds the other two ($t = 4.1$
each) and misses the first ($t = 3.2$). Neither scan finds everything; neither finds anything false. A scan matched to the mechanism, here a window
instead of single lags, has more power, which is why the mechanism should come before the search.

\begin{omfigure}
\begin{tikzpicture}
\begin{axis}[width=10.5cm,height=5.2cm,xlabel={\small lag (days)},ylabel={\small t-statistic, first five years},tick label style={font=\footnotesize},
  xmin=0.5,xmax=10.5,ymin=-5,ymax=5.5,grid=major,grid style={gray!25},legend columns=4,legend style={font=\scriptsize,at={(0.5,-0.3)},anchor=north,
  draw=none,fill=none,/tikz/every even column/.append style={column sep=6pt}},table/col sep=comma]
\addplot[omThm,thick,mark=*] table[x=lag,y=link1]{figdata/strategies-1/22-cross-asset-lead-lag/lagscan.csv};
\addplot[omDef,thick,mark=square*] table[x=lag,y=link2]{figdata/strategies-1/22-cross-asset-lead-lag/lagscan.csv};
\addplot[omExo,thick,mark=triangle*] table[x=lag,y=link3]{figdata/strategies-1/22-cross-asset-lead-lag/lagscan.csv};
\addplot[black!60,dashed] coordinates {(0.5,4.06) (10.5,4.06)};
\addplot[black!60,dashed,forget plot] coordinates {(0.5,-4.06) (10.5,-4.06)};
\legend{link 1,link 2,link 3,Bonferroni (1\,000 tests)}
\end{axis}
\end{tikzpicture}

\omcaption{The lag-by-lag scan on the synthetic industries and commodities, first five years: the t-statistic of each planted link's lagged
correlation at lags of 1 to 10 days, against the Bonferroni threshold for 1\,000 tests. Each link spreads a commodity's return over the next five days.
Data: \texttt{s1\_xasset.daily}.}\label{fig:s1:cross-asset-lead-lag:scan}
\end{omfigure}

The links found are traded in the last five years: each linked industry is held in proportion to its commodity's five-day return in standard deviations,
capped at two, rebalanced daily at 2 basis points per unit traded. The out-of-sample Sharpe ratio is 2.09 before costs and 1.73 after, on two links. That
is high because the planted diffusion is clean and constant; the public record is monthly, weaker, and changes over time.

\section{Intraday lead--lag}

At intraday horizons the leader is usually a futures contract and the followers are the cash instruments that track it (Book~7, chapter~10's
futures-to-cash lead). The chapter simulates twenty sessions of 23\,400 seconds in which a follower sees the leader's efficient price two seconds late,
both observed at random tick times (one a second on average) with noise (\cref{lst:s1:cross-asset-lead-lag:trades}). The Hayashi--Yoshida
cross-correlation of Book~4, chapter~21 handles the asynchronous ticks without a common grid; its peak is at two seconds in every one of the twenty
sessions, and its shape is symmetric about the lead (\cref{fig:s1:cross-asset-lead-lag:ccf}).

\begin{omfigure}
\begin{tikzpicture}
\begin{axis}[width=10cm,height=4.8cm,ybar,bar width=8pt,xlabel={\small lag of the follower (seconds)},ylabel={\small cross-correlation (peak = 1)},
  tick label style={font=\footnotesize},xmin=-6.8,xmax=6.8,ymin=0,ymax=1.1,grid=major,grid style={gray!25},table/col sep=comma]
\addplot[fill=omThm!60,draw=omThm] table[x=lag,y=ccf]{figdata/strategies-1/22-cross-asset-lead-lag/ccf.csv};
\end{axis}
\end{tikzpicture}

\omcaption{The Hayashi--Yoshida cross-correlation between a leader and a follower that sees the leader's price two seconds late, averaged over twenty
simulated sessions and scaled to a peak of one; a peak at a positive lag means the second series follows. Data:
\texttt{s1\_xasset.intraday}.}\label{fig:s1:cross-asset-lead-lag:ccf}
\end{omfigure}

The trade: when the leader has moved more than 2 basis points over the last two seconds, buy or sell the follower in the same direction after a latency,
hold two seconds, and pay half a basis point each way.

\begin{center}
\small
\begin{tabular}{@{}lcccc@{}}
\toprule
latency (seconds) & 0 & 1 & 2 & 3\\
\midrule
trades a session & 2\,343 & 2\,193 & 2\,025 & 1\,862\\
net return per trade (bp) & 0.73 & 0.35 & $-0.51$ & $-0.81$\\
share of trades that gain & 69.6\% & 57.2\% & 32.0\% & 25.0\%\\
sessions with a positive total & 20 & 20 & 0 & 0\\
\bottomrule
\end{tabular}
\end{center}

A trader who acts at once keeps most of the two-second move; one second late keeps half of it; at two seconds the follower has already moved and the
trade pays the spread for nothing. The lead is the same at every latency. What changes is who can use it, as with chapter~17's news.

\section{Slow diffusion and who is slow}

The public record says where slowness comes from. Menzly and Ozbas found that stocks in economically related supplier and customer industries predict
each other's returns, and that the cross-predictability is weaker when more informed investors follow the stocks (more analysts, more institutional
ownership): specialisation makes investors slow to see news in the industries they do not cover. Hou found that the lead of big firms over small ones is
mostly within industries, comes from slow adjustment to negative news, and is strongest in small, less competitive and neglected industries. Hong,
Torous and Valkanov's industries led the market by up to two months, and the industries that did so were those that also forecast economic
activity. Their 2014 note, which extends the sample to 2013 with posted code, finds a robust core of leading industries but a smaller set than before,
with the predictive relations changing over time. Leads decay as they become known, as most published signals do (Book~7, chapter~13).

\section{Trading it}

An intermarket strategy is built in four steps: a mechanism that says who should lead whom and at what speed; a scan matched to it, with the
multiple-testing correction for the number of candidates; a test of stability in time, since leads come and go; and costs, because the predicted
move is small. The daily and intraday versions differ in what limits them. At daily horizons the limit is statistical: a few true links among
hundreds of candidates, each weak. At intraday horizons the limit is speed and cost: the lead is obvious, and it belongs to whoever sees it first
and pays least to trade on it.

\section{Strategy files}

\begin{strategyfile}{Bond-to-equity signal}\label{strat:s1:cross-asset-lead-lag:bond}
\sfield{Who pays you, and why} Equity investors slow to price rate news into rate-sensitive sectors.
\sfield{Instruments and venues} Bond futures as leaders; banks, utilities, real estate and other rate-sensitive stocks or sector futures as followers.
\sfield{Signal} Recent bond futures returns mapped to sectors by their estimated rate sensitivity.
\sfield{Sizing and execution} Sector positions sized by signal and volatility; daily.
\sfield{Costs} Moderate; sector baskets or futures.
\sfield{How it dies} Faster cross-asset traders; changing rate sensitivities.
\sfield{Horizon, capacity, infrastructure} Days.
\sfield{Backtest honestly} Sensitivities estimated from past data only; a scan corrected for every sector and lag tried.
\sfield{Sources} No performance figure verified; the mechanism is \omterm{def:s1:cross-asset-lead-lag:intermarket}{slow diffusion} (Hong, Torous and Valkanov, 2007).
\end{strategyfile}

\begin{strategyfile}{Commodity-to-producer signal}\label{strat:s1:cross-asset-lead-lag:commodity}
\sfield{Who pays you, and why} Equity investors who specialise and follow commodity markets less closely than commodity traders.
\sfield{Instruments and venues} Commodity futures as leaders; producers and heavy users as followers.
\sfield{Signal} The commodity's return over a window matched to the diffusion speed.
\sfield{Sizing and execution} Long producers after commodity rises, short after falls, hedged for the market; daily or weekly.
\sfield{Costs} Moderate.
\sfield{How it dies} More informed investors covering both markets, as Menzly and Ozbas found.
\sfield{Horizon, capacity, infrastructure} Days to weeks.
\sfield{Backtest honestly} The links chosen from a mechanism, then tested; out-of-sample periods.
\sfield{Sources} Hong, Torous and Valkanov (2007); Menzly and Ozbas (2010); this chapter: 1.73 after costs on two planted links.
\end{strategyfile}

\begin{strategyfile}{Futures-to-cash intraday lead}\label{strat:s1:cross-asset-lead-lag:futcash}
\sfield{Who pays you, and why} Liquidity providers in the cash instruments whose quotes update after the futures'.
\sfield{Instruments and venues} Index futures and the ETFs or stocks that track them.
\sfield{Signal} The futures' move over the last seconds.
\sfield{Sizing and execution} Aggressive orders in the follower within the lead; small, many.
\sfield{Costs} Spreads and fees on every trade; the infrastructure for speed.
\sfield{How it dies} Latency competition: one second late keeps half, two seconds late loses.
\sfield{Horizon, capacity, infrastructure} Seconds; co-location and direct feeds.
\sfield{Backtest honestly} Timestamps as received; the follower's quotes, not trades; queue and latency models.
\sfield{Sources} No performance figure verified; this chapter's simulation.
\end{strategyfile}

\begin{strategyfile}{Currency-to-exporter signal}\label{strat:s1:cross-asset-lead-lag:currency}
\sfield{Who pays you, and why} Equity investors slow to price currency moves into exporters' and importers' earnings.
\sfield{Instruments and venues} Currency futures or forwards as leaders; exporters and importers as followers.
\sfield{Signal} Recent currency returns times each firm's estimated foreign revenue share.
\sfield{Sizing and execution} Long exporters after the home currency weakens, short after it strengthens; hedged for the market.
\sfield{Costs} Moderate.
\sfield{How it dies} Firms' hedging; faster analysts.
\sfield{Horizon, capacity, infrastructure} Weeks; segment revenue data.
\sfield{Backtest honestly} Revenue shares as reported at the time.
\sfield{Sources} No performance figure verified.
\end{strategyfile}

\section{Tutorial: same information, different speeds}\label{tut:s1:cross-asset-lead-lag:tut}

\begin{tutorial}
\textbf{Goal.} Plant diffusion from commodities to industries and find it with two scans; trade what is found out of sample; simulate a two-second lead,
find it with the Hayashi--Yoshida estimator and trade it at four latencies. \textbf{End state:} the table and the two figures.
\begin{tutsteps}
\item \textbf{Diffusion and the scans}: planting, the lag-by-lag scan and the Bonferroni detector.
\omcode{code/firm/xasset/firm_xasset.py}{26}{54}{Planted diffusion, the lagged-correlation scan and detection.}\label{lst:s1:cross-asset-lead-lag:scan}
\item \textbf{The intraday trade}: act on the leader's move after a latency.
\omcode{code/firm/xasset/firm_xasset.py}{80}{93}{Trading the follower after a latency.}\label{lst:s1:cross-asset-lead-lag:trades}
\item \textbf{Run} \texttt{daily()} and \texttt{intraday()}, and \texttt{fig\_xasset.py}.
\end{tutsteps}
\textbf{What to change next.} Let a link switch off halfway and test stability with rolling scans; correlate two leaders so that a false link appears
through the correlated one; make the follower's quotes stale only part of the time.
\end{tutorial}

\section{Build: intermarket signals}\label{bld:s1:cross-asset-lead-lag:xasset}

\begin{build}
\bfield{Purpose} Planted diffusion, lead scans with multiple-testing control, \omterm{def:s1:cross-asset-lead-lag:intermarket}{intermarket signals}, and an intraday leader--follower simulator.
\bfield{Interface} \texttt{diffuse(follow, lead, links, beta, lag)}, \texttt{lead\_scan(lead, follow, lags)}, \texttt{detect(tstats, n\_tests, alpha)},
\texttt{signal(lead, links, window, n\_follow)}, \texttt{simulate\_pair(seconds, lead, vol, noise, rate, rng)}, \texttt{lead\_trades(g1, g2, window,
latency, hold, threshold, half\_spread)}.
\bfield{Rules} Signals from past returns only; scans corrected for every test run; one intraday trade at a time.
\bfield{Acceptance tests} \texttt{code/firm/xasset/tests/}: diffusion and signals by hand, a scan that finds exactly one planted lag, a lead trade by
hand.
\bfield{Stretch} Rolling stability tests; false discovery rate control; queue-aware intraday fills.
\end{build}

\begin{omsources}
\item H.~Hong, W.~Torous and R.~Valkanov, ``Do industries lead stock markets?'', \emph{Journal of Financial Economics} 83(2), 2007; and ``Note on Do
Industries Lead Stock Markets'', 2014.
\item L.~Menzly and O.~Ozbas, ``Market segmentation and cross-predictability of returns'', \emph{Journal of Finance} 65(4), 2010.
\item K.~Hou, ``Industry information diffusion and the lead-lag effect in stock returns'', \emph{Review of Financial Studies} 20(4), 2007.
\end{omsources}

\section{Exercises}

\begin{exercise}[$\star$]\label{exo:s1:cross-asset-lead-lag:1}
What two-sided z threshold gives a 5\% family-wise error over 1\,000 tests with Bonferroni's correction? Over 100?
\end{exercise}

\begin{exercise}[$\star$]\label{exo:s1:cross-asset-lead-lag:2}
A follower receives 0.08 times the average of a leader's last five daily returns. With volatilities of 25\% for the leader and 8\% for the follower, what
is the correlation at one lag, and its t-statistic over 1\,260 days? And for the five-day window over 1\,255 days?
\end{exercise}

\begin{exercise}[$\star$]\label{exo:s1:cross-asset-lead-lag:3}
Why does the Hayashi--Yoshida estimator suit tick data that do not arrive at the same times?
\end{exercise}

\begin{exercise}[$\star\star$]\label{exo:s1:cross-asset-lead-lag:4}
The trade at a latency of two seconds loses 0.51 basis points per trade after a cost of one basis point. What does it earn before costs, and why is it
still positive?
\end{exercise}

\begin{exercise}[$\star\star$]\label{exo:s1:cross-asset-lead-lag:5}
Why does the window scan find links the lag-by-lag scan misses, and the reverse?
\end{exercise}

\begin{exercise}[$\star\star$]\label{exo:s1:cross-asset-lead-lag:6}
Two commodities are strongly correlated and only one leads an industry. What can a scan report, and how would you tell?
\end{exercise}

\begin{exercise}[$\star\star\star$]\label{exo:s1:cross-asset-lead-lag:7}
\emph{Coding.} Plant the diffusion at 0.12 instead of 0.08 in \texttt{s1\_xasset.py} and rerun \texttt{daily()}. What does each scan find, and is every
link true?
\end{exercise}

\begin{exercise}[$\star\star\star$]\label{exo:s1:cross-asset-lead-lag:8}
\emph{Find the flaw.} ``We correlated 500 stocks with 60 futures at lags of 1 to 20 days and traded the 50 strongest pairs: backtest Sharpe ratio 3.''
\end{exercise}

\section{Problem: Same Information, Different Speeds}

\begin{problem}[{Weekend problem --- leads by the day and by the second}]\label{pb:s1:cross-asset-lead-lag:1}
The chapter's synthetic markets and the public record.

\medskip\noindent\textbf{Part I --- Diffusion.}
\begin{enumerate}
\item Define an \omterm{def:s1:cross-asset-lead-lag:intermarket}{intermarket signal} and \omterm{def:s1:cross-asset-lead-lag:intermarket}{slow diffusion}.
\item What did Hong, Torous and Valkanov find, and what did their 2014 note add?
\item What did Menzly and Ozbas find about specialisation?
\item What did Hou find about big and small firms?
\end{enumerate}

\medskip\noindent\textbf{Part II --- Daily.}
\begin{enumerate}[resume]
\item Describe the planted diffusion and its strength per lag.
\item Describe the two scans and their thresholds.
\item What does each scan find?
\item What did the found links earn out of sample?
\end{enumerate}

\medskip\noindent\textbf{Part III --- Intraday.}
\begin{enumerate}[resume]
\item Describe the simulated pair and the estimator used.
\item Where is the cross-correlation's peak, and how often?
\item Describe the trade and its costs.
\item Give the net return per trade at each latency.
\end{enumerate}

\medskip\noindent\textbf{Part IV --- The verdict.}
\begin{enumerate}[resume]
\item State the \emph{named result}: the lead detected at each horizon and the strategy's Sharpe ratio after costs.
\item What limits the daily strategy, and what limits the intraday one?
\item Why should the mechanism come before the scan?
\item Why do leads decay?
\item How would you test a lead's stability?
\item Which strategy file needs the fastest infrastructure?
\item How does this chapter relate to chapter~17?
\item In one sentence: who owns a lead?
\end{enumerate}
\end{problem}

\section{Interview questions}

\begin{interviewq}[$\star$ \iqroles{researcher}]\label{iq:s1:cross-asset-lead-lag:1}
How would you test whether one market leads another?
\end{interviewq}

\begin{interviewq}[$\star\star$ \iqroles{researcher}]\label{iq:s1:cross-asset-lead-lag:2}
You scan 10\,000 pairs and find 40 significant at 1\%. What do you conclude?
\end{interviewq}

\begin{interviewq}[$\star\star$ \iqroles{developer}]\label{iq:s1:cross-asset-lead-lag:3}
How would you measure the lead between a futures contract and an ETF from tick data with different timestamps and clocks?
\end{interviewq}

\begin{interviewq}[$\star\star$ \iqroles{trader}]\label{iq:s1:cross-asset-lead-lag:4}
Your futures-to-cash strategy's profit per trade halved in a month. What happened?
\end{interviewq}

\begin{interviewq}[$\star\star$ \iqroles{risk}]\label{iq:s1:cross-asset-lead-lag:5}
What are the risks of an intermarket strategy that holds positions in one market based on another?
\end{interviewq}

\begin{interviewq}[$\star\star\star$ \iqroles{researcher}]\label{iq:s1:cross-asset-lead-lag:6}
A follower's return is $y_t = \beta \frac{1}{L}\sum_{k=1}^{L} x_{t-k} + e_t$ with independent $x$ and $e$. Derive the correlation of $y_t$ with a single lag
$x_{t-k}$ and with the window $\sum_{k=1}^{L} x_{t-k}$, and the ratio of the two t-statistics.
\end{interviewq}
